\section{The SOLID Principles}
\label{sec:lld-solid}

\problemstatement{State each of the five SOLID principles, the design
smell it cures, and a concrete before/after in C++. These principles are
the vocabulary interviewers use to critique a design, so each is presented
as a problem it solves.}

\begin{patternbox}[SOLID in One Line Each]
\textbf{S}ingle Responsibility -- one reason to change per class.
\textbf{O}pen/Closed -- open to extension, closed to modification.
\textbf{L}iskov Substitution -- a subtype must be usable anywhere its base
is. \textbf{I}nterface Segregation -- many small interfaces beat one fat
one. \textbf{D}ependency Inversion -- depend on abstractions, not
concretions.
\end{patternbox}

\subsection*{S -- Single Responsibility Principle}

A class should have exactly one reason to change. An \texttt{Invoice} that
computes totals \emph{and} formats itself as PDF \emph{and} emails itself
has three reasons to change (tax rules, layout, mail server) and three
teams stepping on each other. Split it: \texttt{Invoice} (data + totals),
\texttt{InvoicePrinter} (formatting), \texttt{InvoiceMailer} (sending).

\subsection*{O -- Open/Closed Principle}

Software entities should be open for extension but closed for
modification. The smell is a growing \texttt{switch} on a type tag:

\begin{lstlisting}
// BAD: every new shape edits this function (closed to extension).
double area(const Shape& s) {
    if (s.type == CIRCLE)    return 3.14159 * s.r * s.r;
    if (s.type == RECTANGLE) return s.w * s.h;
    // ... adding Triangle means editing here again
}
\end{lstlisting}

Replace the tag with polymorphism so new shapes are \emph{added}, not
\emph{edited} in:

\begin{lstlisting}
struct Shape { virtual double area() const = 0; virtual ~Shape() = default; };
struct Circle : Shape {
    double r; Circle(double r_) : r(r_) {}
    double area() const override { return 3.14159 * r * r; }
};
struct Rectangle : Shape {
    double w, h; Rectangle(double w_, double h_) : w(w_), h(h_) {}
    double area() const override { return w * h; }
};
// A new Triangle class needs zero changes to existing code.
\end{lstlisting}

\subsection*{L -- Liskov Substitution Principle}

A subclass must honour every promise of its base, so code using the base
never breaks when handed the subclass. The classic violation: a
\texttt{Square} that inherits \texttt{Rectangle} and overrides
\texttt{setWidth} to also change the height silently breaks any code that
sets width and height independently and expects them to stay independent.
The fix is usually to \emph{not} force the inheritance -- model
\texttt{Square} and \texttt{Rectangle} as siblings under a common
\texttt{Shape}, or use composition.

\begin{ideabox}[LSP Is About Behaviour, Not Just Signatures]
Compiling is not enough. A subtype substitutes correctly only if it keeps
the base's \emph{behavioural contract}: it must not strengthen
preconditions, weaken postconditions, or throw ``unsupported'' for methods
the base guarantees. If an override has to disable a base method, the
inheritance is wrong.
\end{ideabox}

\subsection*{I -- Interface Segregation Principle}

No client should be forced to depend on methods it does not use. A fat
\texttt{Machine} interface with \texttt{print()}, \texttt{scan()},
\texttt{fax()} forces a simple printer to stub out \texttt{scan} and
\texttt{fax}. Split into \texttt{Printer}, \texttt{Scanner},
\texttt{Fax}; a multifunction device implements all three, a basic printer
implements only one.

\subsection*{D -- Dependency Inversion Principle}

High-level modules should depend on abstractions, not on concrete
low-level modules -- and the abstraction is injected, not constructed
inside:

\begin{lstlisting}
#include <bits/stdc++.h>
using namespace std;

struct MessageSender {                 // abstraction
    virtual void send(const string& msg) = 0;
    virtual ~MessageSender() = default;
};
struct EmailSender : MessageSender {
    void send(const string& msg) override { cout << "Email: " << msg << "\n"; }
};
struct SmsSender : MessageSender {
    void send(const string& msg) override { cout << "SMS: " << msg << "\n"; }
};

class NotificationService {            // high-level policy
    MessageSender& sender;             // depends on the abstraction
public:
    NotificationService(MessageSender& s) : sender(s) {}   // injected
    void notify(const string& msg) { sender.send(msg); }
};

int main() {
    EmailSender email; SmsSender sms;
    NotificationService viaEmail(email), viaSms(sms);
    viaEmail.notify("Order shipped");
    viaSms.notify("OTP 4321");         // swap implementation, no edits
    return 0;
}
\end{lstlisting}

\begin{takeawaybox}
SOLID is one theme in five costumes: \emph{isolate what changes}. SRP
isolates responsibilities, OCP isolates new behaviour behind polymorphism,
LSP keeps subtypes safe to substitute, ISP isolates clients from
irrelevant methods, and DIP isolates high-level policy from low-level
detail via an injected abstraction. Every pattern in the next chapter is a
reusable structure that upholds one or more of these.
\end{takeawaybox}
